\section{Correspondence Map: This Paper \texorpdfstring{$\leftrightarrow$}{<->} Conceptual Topography}

概念地形論（Conceptual Topography; CT）をあわせて読む読者のために、両論文の主要概念の対応関係をここに示す。

本稿と CT は、同一の差分発展構造の異なる投影窓を記述する。両論文の変数は対応関係にあるが、層をまたいだ単純な同一視ではない。どちらの論文も他を基礎づけはしない。両論文は異なる観測窓から、それぞれ意識として開く状態と、概念として安定する地形を、同一の生成的文法（VED×IFGT）の異なる投影として記述する。本稿 §\ref{sec:projected-uses-stance} で固定された投影的使用の立場は、本対応表全体を貫いている。

\begin{table}[H]
\centering
\small
\begin{tabularx}{\textwidth}{>{\raggedright\arraybackslash}X >{\raggedright\arraybackslash}X >{\raggedright\arraybackslash}X}
\toprule
本稿 & CT（Conceptual Topography） & 共有変数・対応関係 \\
\midrule
低速層の沈殿構造 & 情報ポテンシャル$\Phi$の地形 & 両者は$\Phi^*$として収束する長時間アトラクタ \\
SSO（時間的非閉包） & Invisible GPU（意識下で動く地形） & $\Phi$が意識に直接参照されず思考を制約する \\
TCA（速度差の変換） & トークン結合$\lambda\sum_i \ell_i \cdot \nabla\Phi$ & 速度差／空間差を変換する同型の構造 \\
意味勾配$\nabla\Phi$によるログ参照方向 & 概念定義$\nabla\Phi\approx 0$の局所安定領域 & $\nabla\Phi$は意識の偏りと概念の形成を同時に規定する \\
ログ参照（意識が開かれる状態） & 通過コスト低下（理解が成立する状態） & 両者とも$\nabla\Phi$構造への参照・通過として記述される \\
クオリア（沈殿境界のライブな前景） & インサイト（急激な位相転移） & 同一の境界現象を異なる時間粒度で観察したもの \\
$I(x,\tau) > 0$（完全閉包の排除） & $I(x,\tau) > 0$（完全概念閉包の排除） & 正値性は各観測窓で通常状態としての完全閉包が排除されることを示すが、意識と概念地形を直同一視するものではない \\
\bottomrule
\end{tabularx}
\end{table}

\textbf{読み方の提案：}

本稿を先に読んだ場合、CT を読むことで、本稿の「意識が開かれる時間条件」が「概念が形成される空間条件」と同一の構造から来ていることが見えてくる。意識状態の中で何が参照されるかは、その人が持つ概念地形の幾何学によって決まる。

CT を先に読んだ場合、本稿を読むことで、CT の「思考が地形を流れる」という記述の背後に、高速層と低速層の速度差が必要であることが見えてくる。地形は所与ではなく、SSO が成立しているときにのみ参照可能になる。

相互参照のもとで読むとき、両論文は次のことを示唆する。意識として開く状態と、概念として安定する地形は、同一の VED×IFGT 文法の異なる投影安定化レジームであって、同一の実体でも還元可能な対象でもない。同じ方針は VED 物理層・Bio-IFGT 生命層・知性編知性層 \cite{Intelligence} にも拡張され、四層全体が同一の差分発展構造の異なる投影窓として読まれる。本稿の意識窓はそのうちのひとつであり、他の窓との対応は還元ではなく、各窓に固有の有効座標のもとでの構造的同型性として理解される。
